\chapter{2011年安徽卷（理）}

\section{单选题}

\begin{problem}
设 \(\i\) 是虚数单位，复数 \(\frac{1 + a\i}{2 - \i}\) 为纯虚数，则实数 \(a\) 为（\quad）

\choices
    {\(2\)}
    {\(-2\)}
    {\(-\frac{1}{2}\)}
    {\(\frac{1}{2}\)}
\end{problem}

\begin{answer}
A
\end{answer}

\begin{solution}
分母乘以共轭复数，得
\[
\frac{1+a\i}{2-\i}=\frac{(1+a\i)(2+\i)}{5}=\frac{2-a+(1+2a)\i}{5}
\]
复数为纯虚数，所以实部为零，即 \(2-a=0\)，解得 \(a=2\)，故选 A.
\end{solution}

\begin{problem}
双曲线 \(2x^{2} - y^{2} = 8\) 的实轴长是（\quad）

\choices
    {\(2\)}
    {\(2\sqrt{2}\)}
    {\(4\)}
    {\(4\sqrt{2}\)}
\end{problem}

\begin{answer}
C
\end{answer}

\begin{solution}
将方程化为标准形式：
\[
\frac{x^{2}}{4}-\frac{y^{2}}{8}=1
\]
因此双曲线的实半轴长为 \(2\)，实轴长为 \(2\times2=4\)，故选 C.
\end{solution}

\begin{problem}
设 \( f(x) \) 是定义在 \( \R \) 上的奇函数，当 \( x \leqslant 0 \) 时，\( f(x) = 2x^2 - x \)，则 \( f(1) = \)（\quad）

\choices
    {\(-3\)}
    {\(-1\)}
    {\(1\)}
    {\(3\)}
\end{problem}

\begin{answer}
A
\end{answer}

\begin{solution}
因为 \(-1\leqslant0\)，所以
\[
f(-1)=2(-1)^2-(-1)=3
\]
函数为奇函数，故 \(f(1)=-f(-1)=-3\)，故选 A.
\end{solution}

\begin{problem}
设变量 \(x\)，\(y\) 满足 \(|x| + |y| \leqslant 1\)，则 \(x + 2y\) 的最大值和最小值分别为（\quad）

\choices
    {\(1\)，\(-1\)}
    {\(2\)，\(-2\)}
    {\(1\)，\(-2\)}
    {\(2\)，\(-1\)}
\end{problem}

\begin{answer}
B
\end{answer}

\begin{solution}
由三角不等式，
\[
x+2y\leqslant |x|+2|y|=(|x|+|y|)+|y|\leqslant2
\]
当 \((x,y)=(0,1)\) 时取到上界. 类似地，
\[
x+2y\geqslant-|x|-2|y|\geqslant-2
\]
当 \((x,y)=(0,-1)\) 时取到下界. 因此最大值和最小值分别为 \(2\)，\(-2\)，故选 B.
\end{solution}

\begin{problem}
在极坐标系中，点 \(\left(2, \frac{\pi}{3}\right)\) 到圆 \(\rho = 2\cos \theta\) 的圆心的距离为（\quad）

\choices
    {\(2\)}
    {\(\sqrt{4 + \frac{\pi^2}{9}}\)}
    {\(\sqrt{1 + \frac{\pi^2}{9}}\)}
    {\(\sqrt{3}\)}
\end{problem}

\begin{answer}
D
\end{answer}

\begin{solution}
极坐标点 \(\left(2,\frac{\pi}{3}\right)\) 的直角坐标为
\[
\left(2\cos\frac{\pi}{3},2\sin\frac{\pi}{3}\right)=(1,\sqrt3)
\]
圆的方程满足 \(\rho^2=2\rho\cos\theta\)，所以其直角坐标方程为 \(x^2+y^2=2x\)，即 \((x-1)^2+y^2=1\)，圆心为 \((1,0)\). 所求距离为
\[
\sqrt{(1-1)^2+\left(\sqrt3-0\right)^2}=\sqrt3
\]
故选 D.
\end{solution}

\begin{problem}
一个空间几何体的三视图如图所示，则该几何体的表面积为（\quad）

\bitmapfigure[max width=0.62\linewidth]{img/2011/anhui_science/q6_fig1.png}

\choices
    {\(48\)}
    {\(32 + 8\sqrt{17}\)}
    {\(48 + 8\sqrt{17}\)}
    {\(80\)}
\end{problem}

\begin{answer}
C
\end{answer}

\begin{solution}
由三视图可知，该几何体是以等腰梯形为底面的直棱柱. 等腰梯形的上底、下底和高分别为 \(2\)，\(4\)，\(4\)，直棱柱的高为 \(4\). 梯形两腰长为
\[
\sqrt{4^2+\left(\frac{4-2}{2}\right)^2}=\sqrt{17}
\]
因此一个底面的面积为
\[
\frac{2+4}{2}\times4=12
\]
两个底面的面积和为 \(24\)，四个侧面的面积和为
\[
4\left(2+4+2\sqrt{17}\right)=24+8\sqrt{17}
\]
所以表面积为 \(24+24+8\sqrt{17}=48+8\sqrt{17}\)，故选 C.
\end{solution}

\begin{problem}
命题“所有能被 \(2\) 整除的整数都是偶数”的否定是（\quad）

\choices
    {所有不能被 \(2\) 整除的整数都是偶数}
    {所有能被 \(2\) 整除的整数都不是偶数}
    {存在一个不能被 \(2\) 整除的整数是偶数}
    {存在一个能被 \(2\) 整除的整数不是偶数}
\end{problem}

\begin{answer}
D
\end{answer}

\begin{solution}
原命题是“对所有整数，若该整数能被 \(2\) 整除，则它是偶数”. 否定全称命题时，将“所有”改为“存在一个”，并否定结论，因此否定为“存在一个能被 \(2\) 整除的整数不是偶数”，故选 D.
\end{solution}

\begin{problem}
设集合 \(A = \{1, 2, 3, 4, 5, 6\}\)，\(B = \{4, 5, 6, 7, 8\}\)，则满足 \(S \subseteq A\) 且 \(S \cap B \neq \varnothing\) 的集合 \(S\) 的个数是（\quad）

\choices
    {\(57\)}
    {\(56\)}
    {\(49\)}
    {\(8\)}
\end{problem}

\begin{answer}
B
\end{answer}

\begin{solution}
集合 \(A\) 有 \(2^6=64\) 个子集. 因为 \(S\subseteq A\)，所以 \(S\cap B=S\cap\{4,5,6\}\). 不满足条件的子集必须不含 \(4\)，\(5\)，\(6\)，只能从 \(\{1,2,3\}\) 中取，计有 \(2^3=8\) 个. 因此满足条件的集合数为
\[
2^6-2^3=64-8=56
\]
故选 B.
\end{solution}

\begin{problem}
已知函数 \( f(x) = \sin (2x + \varphi) \)，其中 \( \varphi \) 为实数，若 \( f(x) \leqslant \left|f\left(\frac{\pi}{6}\right)\right| \) 对 \( x \in \R \) 恒成立，且 \( f\left(\frac{\pi}{2}\right) > f(\pi) \)，则 \( f(x) \) 的单调递增区间是（\quad）

\choices
    {\(\left[k\pi - \frac{\pi}{3}, k\pi + \frac{\pi}{6}\right]\)（\(k \in \Z\)）}
    {\(\left[k\pi, k\pi + \frac{\pi}{2}\right]\)（\(k \in \Z\)）}
    {\(\left[k\pi + \frac{\pi}{6}, k\pi + \frac{2\pi}{3}\right]\)（\(k \in \Z\)）}
    {\(\left[k\pi - \frac{\pi}{2}, k\pi\right]\)（\(k \in \Z\)）}
\end{problem}

\begin{answer}
C
\end{answer}

\begin{solution}
函数 \(f\) 的最大值为 \(1\)，由题设恒成立不等式可知
\[
1\leqslant\left|f\left(\frac{\pi}{6}\right)\right|\leqslant1
\]
故 \(\left|\sin\left(\frac{\pi}{3}+\varphi\right)\right|=1\)，从而
\(\varphi=\frac{\pi}{6}+k\pi\)，\(k\in\Z\).
又
\(f\left(\frac{\pi}{2}\right)=-\sin\varphi\)，\(qquad f(\pi)=\sin\varphi\)
由题设得 \(\sin\varphi<0\). 于是 \(k\) 为奇数，故可将 \(\varphi\) 写成 \(\frac{7\pi}{6}+2m\pi\).

因为
\[
f'(x)=2\cos(2x+\varphi)
\]
所以函数递增当且仅当 \(-\frac{\pi}{2}+2j\pi\leqslant2x+\varphi\leqslant\frac{\pi}{2}+2j\pi\). 代入 \(\varphi=\frac{7\pi}{6}+2m\pi\)，并整理得
\[
j\pi-\frac{5\pi}{6}\leqslant x\leqslant j\pi-\frac{\pi}{3}
\]
将整数下标重新记为 \(k-1\)，递增区间为 \(\left[k\pi+\frac{\pi}{6},k\pi+\frac{2\pi}{3}\right]\)，故选 C.
\end{solution}

\begin{problem}
函数 \( f(x) = ax^{m}(1 - x)^n \) 在区间 \([0,1]\) 上的图象如图所示，则 \(m\)，\(n\) 的值可能是（\quad）

\bitmapfigure[max width=0.46\linewidth]{img/2011/anhui_science/q10_fig1.png}

\choices
    {\(m = 1, n = 1\)}
    {\(m = 1\)，\(n = 2\)}
    {\(m = 2\)，\(n = 1\)}
    {\(m = 3\)，\(n = 1\)}
\end{problem}

\begin{answer}
B
\end{answer}

\begin{solution}
对候选的正整数 \(m\)，\(n\)，有
\[
f'(x)=a x^{m-1}(1-x)^{n-1}\bigl[m-(m+n)x\bigr]
\]
图象的最高点位于 \(0.5\) 左侧，因此极大值点应满足 \(\frac{m}{m+n}<\frac12\)，即 \(m<n\). 四个选项中只有 \(m=1\)，\(n=2\) 满足这一条件；并且此时右端点因子为 \((1-x)^2\)，图象在 \(x=1\) 处与横轴相切，与图形一致，故选 B.
\end{solution}

\section{填空题}

\begin{problem}
如图所示，程序框图（算法流程图）的输出结果是\fillinblank{}.

\bitmapfigure[max width=0.42\linewidth]{img/2011/anhui_science/q11_fig1.png}
\end{problem}

\begin{answer}
\(15\)
\end{answer}

\begin{solution}
初始时 \(T=0\)，\(k=0\). 每次先执行 \(T=T+k\)，若结果不超过 \(105\)，再令 \(k=k+1\) 并进入下一轮. 因此依次累加 \(0\)，\(1\)，\(\ldots\)，\(14\) 后，
\[
T=0+1+\cdots+14=\frac{14\times15}{2}=105
\]
此时条件 \(T>105\) 不成立，令 \(k=15\). 下一轮得到 \(T=105+15=120>105\)，所以输出 \(k=15\).
\end{solution}

\begin{problem}
设 \((x - 1)^{21} = a_0 + a_1x + a_2x^2 + \dots + a_{21}x^{21}\)，则 \(a_{10} + a_{11} =\fillinblank{}\).
\end{problem}

\begin{answer}
\(0\)
\end{answer}

\begin{solution}
由二项式定理，\(x^j\) 的系数为 \(a_j=\mathrm C_{21}^{j}(-1)^{21-j}\). 于是
\(a_{10}=\mathrm C_{21}^{10}(-1)^{11}=-\mathrm C_{21}^{10}\)，\(a_{11}=\mathrm C_{21}^{11}(-1)^{10}=\mathrm C_{21}^{11}\).
又 \(\mathrm C_{21}^{10}=\mathrm C_{21}^{11}\)，故 \(a_{10}+a_{11}=0\).
\end{solution}

\begin{problem}
已知向量 \(\bs{a}\)，\(\bs{b}\) 满足 \((\bs{a} + 2\bs{b}) \cdot (\bs{a} - \bs{b}) = -6\)，且 \(|\bs{a}| = 1\)，\(|\bs{b}| = 2\)，则 \(\bs{a}\) 与 \(\bs{b}\) 的夹角为\fillinblank{}.
\end{problem}

\begin{answer}
\(\frac{\pi}{3}\)
\end{answer}

\begin{solution}
展开数量积，得
\[
(\bs{a}+2\bs{b})\cdot(\bs{a}-\bs{b})=|\bs{a}|^2+\bs{a}\cdot\bs{b}-2|\bs{b}|^2
\]
代入已知条件，
\[
1+\bs{a}\cdot\bs{b}-2\times4=-6
\]
所以 \(\bs{a}\cdot\bs{b}=1\). 设两向量夹角为 \(\theta\)，则
\[
\cos\theta=\frac{\bs{a}\cdot\bs{b}}{|\bs{a}|\,|\bs{b}|}=\frac12
\]
由于 \(0\leqslant\theta\leqslant\pi\)，故 \(\theta=\frac{\pi}{3}\).
\end{solution}

\begin{problem}
已知 \(\triangle ABC\) 的一个内角为 \(120^{\circ}\)，并且三边长构成公差为 \(4\) 的等差数列，则 \(\triangle ABC\) 的面积为\fillinblank{}.
\end{problem}

\begin{answer}
\(15\sqrt{3}\)
\end{answer}

\begin{solution}
设三边按从小到大排列为 \(x-4\)，\(x\)，\(x+4\)，其中最大边 \(x+4\) 对应最大角 \(120^{\circ}\). 由余弦定理，
\[
(x+4)^2=x^2+(x-4)^2-2x(x-4)\cos120^{\circ}=x^2+(x-4)^2+x(x-4)
\]
化简得 \(2x(x-10)=0\). 因边长为正且 \(x>4\)，所以 \(x=10\)，三边长为 \(6\)，\(10\)，\(14\). 因此面积为
\[
\frac12\times6\times10\times\sin120^{\circ}=15\sqrt3
\]
\end{solution}

\begin{problem}
在平面直角坐标系中，如果 \(x\) 与 \(y\) 都是整数，就称点 \((x, y)\) 为整点，下列命题中正确的是\fillinblank{}. （写出所有正确命题的编号）

\begin{circlelist}
\item 存在这样的直线，既不与坐标轴平行又不经过任何整点；
\item 如果 \(k\) 与 \(b\) 都是无理数，则直线 \(y = kx + b\) 不经过任何整点；
\item 直线 \(l\) 经过无穷多个整点，当且仅当 \(l\) 经过两个不同的整点；
\item 直线 \(y = kx + b\) 经过无穷多个整点的充分必要条件是：\(k\) 与 \(b\) 都是有理数；
\item 存在恰经过一个整点的直线.
\end{circlelist}
\end{problem}

\begin{answer}
\(1\)，\(3\)，\(5\)
\end{answer}

\begin{solution}
对命题 1，取直线 \(y=\sqrt2x+\frac12\). 它的斜率非零且有限，故不与坐标轴平行；当 \(x\) 为非零整数时，\(\sqrt2x+\frac12\) 为无理数，当 \(x=0\) 时纵坐标为 \(\frac12\)，所以该直线不经过任何整点，命题 1 正确.

命题 2 不正确. 例如取 \(k=\sqrt2\)，\(b=1-\sqrt2\)，二者均为无理数，但直线经过整点 \((1,1)\).

对命题 3，若直线经过无穷多个整点，必然经过其中两个不同的整点. 反之，若直线经过两个不同的整点 \((x_1,y_1)\)，\((x_2,y_2)\)，令 \(d=(x_2-x_1,y_2-y_1)\)，则对任意整数 \(t\)，点 \((x_1,y_1)+td\) 都是整点且在该直线上，并且随 \(t\) 取不同整数而互不相同，所以直线经过无穷多个整点，命题 3 正确.

命题 4 不正确. 例如直线 \(y=\frac12x+\frac13\) 中的斜率和截距都是有理数，但若 \(x\)，\(y\) 都是整数，则 \(6y=3x+2\). 左边模 \(3\) 为 \(0\)，右边模 \(3\) 为 \(2\)，不可能相等，因此该直线没有整点，更不可能经过无穷多个整点.

对命题 5，直线 \(y=\sqrt2x\) 经过整点 \((0,0)\). 若 \(x\) 是非零整数，则 \(y=\sqrt2x\) 是无理数，不可能为整数，因此它恰经过一个整点，命题 5 正确. 故正确命题的编号为 \(1\)，\(3\)，\(5\).
\end{solution}

\section{解答题}

\begin{problem}
设 \( f(x) = \frac{\e^x}{1 + ax^2} \)，其中 \( a \) 为正实数.

\begin{enumerate}
\item 当 \(a=\frac{4}{3}\) 时，求 \( f(x) \) 的极值点；
\item 若 \( f(x) \) 为 \(\R\) 上的单调函数，求 \( a \) 的取值范围.
\end{enumerate}
\end{problem}

\begin{solution}
\begin{enumerate}
\item 对任意 \(a>0\)，分母 \(1+ax^2>0\)，且
\[
f'(x)=\frac{\e^x\bigl(ax^2-2ax+1\bigr)}{(1+ax^2)^2}
\]
当 \(a=\frac43\) 时，
\[
ax^2-2ax+1=\frac13(2x-1)(2x-3)
\]
所以 \(f'(x)\) 在 \(\left(-\infty,\frac12\right)\) 和 \(\left(\frac32,+\infty\right)\) 上为正，在 \(\left(\frac12,\frac32\right)\) 上为负. 因此极值点分别为 \(\left(\frac12,\frac{3\sqrt{\e}}4\right)\) 和 \(\left(\frac32,\frac{\e^{3/2}}4\right)\)，前者为极大值点，后者为极小值点.
\item 因为 \(\e^x\) 和分母的平方始终为正，\(f'(x)\) 的符号由
\[
q(x)=ax^2-2ax+1=a(x-1)^2+1-a
\]
决定. 当 \(0<a<1\) 时，\(q(x)>0\)；当 \(a=1\) 时，\(q(x)=(x-1)^2\geqslant0\)，故此时 \(f\) 在 \(\R\) 上单调递增. 当 \(a>1\) 时，\(q(x)=0\) 有两个根 \(1-\sqrt{1-\frac1a}\) 和 \(1+\sqrt{1-\frac1a}\)，且 \(q(x)\) 在两根之间为负、两根外为正，函数先增、后减、再增，不是单调函数. 因此 \(a\) 的取值范围为 \(0<a\leqslant1\).
\end{enumerate}
\end{solution}

\begin{problem}
如图，\(ABEDFC\) 为多面体，平面 \(ABED\) 与平面 \(ACFD\) 垂直，点 \(O\) 在线段 \(AD\) 上，\(OA = 1\)，\(OD = 2\)，\(\triangle OAB\)，\(\triangle OAC\)，\(\triangle ODE\)，\(\triangle ODF\) 都是正三角形.

\begin{enumerate}
\item 证明直线 \(BC \parallel EF\)；
\item 求棱锥 \(F-OBED\) 的体积.
\end{enumerate}

\bitmapfigure[max width=0.44\linewidth]{img/2011/anhui_science/q17_fig1.png}
\end{problem}

\begin{solution}
\begin{enumerate}
\item 建立直角坐标系，使 \(A=(0,0,0)\)，\(AD\) 为 \(x\) 轴，平面 \(ABED\) 为 \(xOy\) 平面，平面 \(ACFD\) 为 \(xOz\) 平面. 由 \(OA=1\)，\(OD=2\) 得
\(O=(1,0,0)\)，\(D=(3,0,0)\).
按图取各正三角形所在的一侧，则
\(B=\left(\frac12,-\frac{\sqrt3}{2},0\right)\)，\(C=\left(\frac12,0,\frac{\sqrt3}{2}\right)\)
\(E=(2,-\sqrt3,0)\)，\(F=(2,0,\sqrt3)\).
于是
\(\overrightarrow{BC}=\left(0,\frac{\sqrt3}{2},\frac{\sqrt3}{2}\right)\)，\(\overrightarrow{EF}=(0,\sqrt3,\sqrt3)=2\overrightarrow{BC}\)
故 \(BC\parallel EF\).
\item 底面 \(OBED\) 在平面 \(ABED\) 内. 沿对角线 \(OE\) 分成两个三角形，由坐标或底乘高公式，
\(S_{\triangle OBE}=\frac12\left|\left(-\frac12\right)(-\sqrt3)-\left(-\frac{\sqrt3}{2}\right)(1)\right|=\frac{\sqrt3}{2}\)，\(S_{\triangle OED}=\frac12\times OD\times\sqrt3=\sqrt3\).
所以底面积为 \(S_{OBED}=\frac{3\sqrt3}{2}\). 点 \(F\) 到底面所在平面的距离为 \(\sqrt3\)，故棱锥体积为
\[
V_{F-OBED}=\frac13S_{OBED}\cdot\sqrt3=\frac13\times\frac{3\sqrt3}{2}\times\sqrt3=\frac32
\]
\end{enumerate}
\end{solution}

\begin{problem}
在数 \(1\) 和 \(100\) 之间插入 \(n\) 个实数，使得这 \(n + 2\) 个数构成递增的等比数列，将这 \(n + 2\) 个数的乘积记作 \(T_{n}\)，再令 \(a_{n} = \lg T_{n}\)，\(n \geqslant 1\).

\begin{enumerate}
\item 求数列 \(\{a_{n}\}\) 的通项公式；
\item 设 \(b_{n} = \tan a_{n} \cdot \tan a_{n+1}\)，求数列 \(\{b_{n}\}\) 的前 \(n\) 项和 \(S_{n}\).
\end{enumerate}
\end{problem}

\begin{solution}
\begin{enumerate}
\item 设等比数列的公比为 \(r\). 因为数列递增且首末两项为 \(1,100\)，所以 \(r>1\)，并且 \(r^{n+1}=100\)，从而 \(\lg r=\frac{2}{n+1}\). 这 \(n+2\) 项为 \(1\)，\(r\)，\(r^2\)，\(\ldots\)，\(r^{n+1}\)，故
\[
T_n=r^{0+1+\cdots+(n+1)}=r^{\frac{(n+1)(n+2)}2}
\]
因此
\[
a_n=\lg T_n=\frac{(n+1)(n+2)}2\lg r=n+2
\]
\item 由上式，\(b_n=\tan(n+2)\tan(n+3)\). 利用正切差角公式
\[
\tan1=\frac{\tan(j+1)-\tan j}{1+\tan(j+1)\tan j}
\]
得
\[
\tan j\tan(j+1)=\frac{\tan(j+1)-\tan j}{\tan1}-1
\]
令 \(j=3\)，\(4\)，\(\ldots\)，\(n+2\) 并求和，得到
\[
S_n=\sum_{j=3}^{n+2}\tan j\tan(j+1)
=\frac{\tan(n+3)-\tan3}{\tan1}-n
\]
\end{enumerate}
\end{solution}

\begin{problem}
\begin{enumerate}
\item 设 \(x \geqslant 1\)，\(y \geqslant 1\)，证明 \(x + y + \frac{1}{xy} \leqslant \frac{1}{x} + \frac{1}{y} + xy\)；
\item \(1 < a \leqslant b \leqslant c\)，证明 \(\log_{a} b + \log_{b} c + \log_{c} a \leqslant \log_{b} a + \log_{c} b + \log_{a} c\).
\end{enumerate}
\end{problem}

\begin{solution}
\begin{enumerate}
\item 因为 \(xy>0\)，原不等式等价于
\[
x^2y^2-x^2y-xy^2+x+y-1\geqslant0
\]
左端可因式分解为
\[
(x-1)(y-1)(xy-1)
\]
由 \(x\geqslant1\)，\(y\geqslant1\) 可知三个因式均非负，所以原不等式成立.
\item 令 \(u=\ln a\)，\(v=\ln b\)，\(w=\ln c\)，则 \(0<u\leqslant v\leqslant w\). 由换底公式，所要证明的不等式等价于
\[
\frac vu+\frac wv+\frac uw\leqslant\frac uv+\frac vw+\frac wu
\]
再令 \(x=\frac vu\geqslant1\)，\(y=\frac wv\geqslant1\)，则 \(\frac wu=xy\)，上式化为
\[
x+y+\frac1{xy}\leqslant\frac1x+\frac1y+xy
\]
这正是第（1）问已经证明的不等式，故原不等式成立.
\end{enumerate}
\end{solution}

\begin{problem}
工作人员需进入核电站完成某项具有高辐射危险的任务，每次只派一个人进去，且每个人只派一次，工作时间不超过 \(10\) 分钟，如果前一个人 \(10\) 分钟内不能完成任务则撤出，再派下一个人. 现在一共只有甲，乙，丙三个人可派. 他们各自能完成任务的概率分别 \(p_{1}\)，\(p_{2}\)，\(p_{3}\). 假设 \(p_{1}\)，\(p_{2}\)，\(p_{3}\) 互不相等，且假定各人能否完成任务的事件相互独立.

\begin{enumerate}
\item 如果按甲最先，乙次之，丙最后的顺序派人，求任务能被完成的概率. 若改变三个人被派出的先后顺序，任务能被完成的概率是否发生变化？
\item 若按某指定顺序派人，这三个人各自能完成任务的概率依次为 \(q_{1}\)，\(q_{2}\)，\(q_{3}\)，其中 \(q_{1}\)，\(q_{2}\)，\(q_{3}\) 是 \(p_{1}\)，\(p_{2}\)，\(p_{3}\) 的一个排列. 求所需派出人员数目 \(X\) 的分布列和均值（数学期望）\(E(X)\)；
\item 假定 \(1 > p_{1} > p_{2} > p_{3}\)，试分析以怎样的先后顺序派出人员，可使所需派出的人员数目的均值（数学期望）达到最小.
\end{enumerate}
\end{problem}

\begin{solution}
\begin{enumerate}
\item 任务不能完成，当且仅当甲、乙、丙都没有完成任务. 由相互独立性，其概率为 \((1-p_1)(1-p_2)(1-p_3)\)，所以任务能完成的概率为
\[
1-(1-p_1)(1-p_2)(1-p_3)
\]
这个事件只由三个人是否成功决定，与派出顺序无关，因此改变顺序不会改变任务能被完成的概率.
\item 当第一个人完成任务时，\(X=1\)，所以 \(P(X=1)=q_1\)；当第一个人失败而第二个人成功时，\(X=2\)，所以 \(P(X=2)=(1-q_1)q_2\)；当前两个人都失败时，必须派出第三个人，无论第三个人成功与否均有 \(X=3\)，所以 \(P(X=3)=(1-q_1)(1-q_2)\). 这三个概率之和为 \(1\)，故分布列由上述三式给出. 数学期望为
\[
E(X)=q_1+2(1-q_1)q_2+3(1-q_1)(1-q_2)=3-2q_1-q_2+q_1q_2
\]
\item 先比较前两人的顺序. 若交换前两人的成功概率，期望之差为
\[
E(q_1,q_2)-E(q_2,q_1)=q_2-q_1
\]
因此前两人的成功概率中较大的应放在前面. 再固定第一人的概率为 \(q_1<1\)，比较第二、第三人的顺序，有
\[
E(q_1,q_2)-E(q_1,q_3)=(q_3-q_2)(1-q_1)
\]
若 \(q_2>q_3\)，上式小于零，所以第二人应选成功概率较大的. 由 \(p_1>p_2>p_3\)，先甲、再乙、后丙时均值最小.
\end{enumerate}
\end{solution}

\begin{problem}
设 \(\lambda > 0\)，点 \(A\) 的坐标为 \((1,1)\)，点 \(B\) 在抛物线 \(y = x^2\) 上运动，点 \(Q\) 满足 \(\overrightarrow{BQ} = \lambda \overrightarrow{QA}\)，经过点 \(Q\) 与 \(x\) 轴垂直的直线交抛物线于点 \(M\)，点 \(P\) 满足 \(\overrightarrow{QM} = \lambda \overrightarrow{MP}\)，求点 \(P\) 的轨迹方程.

\bitmapfigure[max width=0.42\linewidth]{img/2011/anhui_science/q21_fig1.png}
\end{problem}

\begin{solution}
设 \(B=(t,t^2)\)，并设 \(Q=(u,v)\). 由 \(\overrightarrow{BQ}=\lambda\overrightarrow{QA}\)，有
\[
(u-t,v-t^2)=\lambda(1-u,1-v)
\]
因此
\(u=\frac{t+\lambda}{1+\lambda}\)，\(v=\frac{t^2+\lambda}{1+\lambda}\).
过 \(Q\) 的竖直直线与抛物线交于 \(M=(u,u^2)\). 再由 \(\overrightarrow{QM}=\lambda\overrightarrow{MP}\)，设 \(P=(u,w)\)，则
\(u^2-v=\lambda(w-u^2)\)，\(w=u^2+\frac{u^2-v}{\lambda}\).
由 \(t=(1+\lambda)u-\lambda\)，可得
\[
v=\frac{\bigl((1+\lambda)u-\lambda\bigr)^2+\lambda}{1+\lambda}
=(1+\lambda)u^2-2\lambda u+\lambda
\]
所以
\(u^2-v=-\lambda(u-1)^2\)，\(w=u^2-(u-1)^2=2u-1\).
点 \(P\) 的横坐标 \(u=\frac{t+\lambda}{1+\lambda}\) 随 \(t\) 取遍全体实数，故轨迹为整条直线
\[
y=2x-1
\]
\end{solution}
